%% Journal draft: how accurate must a digital twin be for fault diagnosis?
%%
%% Status: first draft, written without the authors' input on the instruction
%% to proceed on best judgement. Every number is generated by the scripts in
%% ../ and recorded in ../_results.json; ../claims.py checks the text against
%% them. Decisions taken without the authors are listed in ../LABNOTEBOOK.md.
%%
%% Build:  pdflatex paper && bibtex paper && pdflatex paper && pdflatex paper

\documentclass[preprint,12pt]{elsarticle}

\usepackage{amsmath,amssymb,amsthm}
\usepackage{graphicx}
\usepackage{booktabs}
\usepackage{url}
\usepackage[T1]{fontenc}
\usepackage{xcolor}
\usepackage{hyperref}
\hypersetup{colorlinks=true,linkcolor=blue,citecolor=blue,urlcolor=blue}
\graphicspath{{figures/}}

\newtheorem{theorem}{Theorem}
\newtheorem{corollary}{Corollary}
\newtheorem{proposition}[theorem]{Proposition}

\journal{[to be decided by the authors]}

\begin{document}

\begin{frontmatter}

\title{Validated for prediction, unverified for diagnosis: how accurate a
digital twin must be to attribute process faults in laser powder bed fusion}

\author[a]{Soham Jain\corref{cor1}}
\ead{f20240343@pilani.bits-pilani.ac.in}
\cortext[cor1]{Corresponding author.}
\author[a]{Tribeni Roy}
\affiliation[a]{organization={Department of Mechanical Engineering,
                Birla Institute of Technology and Science Pilani},
                city={Pilani}, country={India}}

\begin{abstract}
%%ABSTRACT%%
\end{abstract}

\begin{keyword}
digital twin \sep model validation \sep fault diagnosis \sep identifiability
\sep laser powder bed fusion \sep twin experiment
\end{keyword}

\end{frontmatter}

%% ===========================================================================
\section{Introduction}
\label{sec:intro}

%%INTRO%%

%% ===========================================================================
\section{Setting}
\label{sec:setting}

\subsection{The twin and its two uses}
\label{sec:twin}

The twin is the thermal model of the two companion papers: an Eagar--Tsai
moving point source with an analytic keyhole depth correction, calibrated on
seven single tracks of the 2022 Additive Manufacturing Benchmark on nickel alloy
718 \cite{Weaver2024}. Four parameters are fitted: effective absorptivity
$\mu$, conductivity $k$, heat capacity $c_p$ and a keyhole coefficient $C$.
The calibrated values are $\mu = 0.384$, $k = 35.0$\,W\,m$^{-1}$\,K$^{-1}$,
$c_p = 450$\,J\,kg$^{-1}$\,K$^{-1}$ and $C = 1.00$, and the model reproduces the
measured widths and depths with a mean absolute error of 12.4\,\%.

The twin is used in two ways. For \emph{prediction} it returns the melt pool
width and depth of a healthy build at a given power and speed. This is the
output that calibration fits and that validation checks. For \emph{diagnosis}
it returns, for each of five process faults, the direction in which that fault
moves the readings of three coaxial instruments: a photodiode, a camera
reporting melt pool area, and a two-colour pyrometer. The five faults are a
laser power deviation, a scan speed error, a powder layer thickness deviation, a
degradation of shielding gas flow, and local heat accumulation caused by part
geometry, each entering the model through a stated physical route. Their
whitened signatures, the columns of the relative-sensitivity Jacobian divided by
each instrument's noise, are the input to the diagnostic test of the
certificate paper.

\subsection{The certificate}
\label{sec:certificate}

Let $z_w \in \mathbb{R}^S$ be an observed change in the $S$ instrument
readings, expressed in units of each instrument's noise, and let $u_k$ be the
unit vector along the whitened signature of fault $k$. The fault magnitude is
unknown and is removed by least squares, which leaves for each candidate the
profile residual
\begin{equation}
R_k = \|z_w\|^2 - \langle z_w, u_k\rangle^2 .
\label{eq:residual}
\end{equation}
If fault $k$ is the true cause and the observation is linear and Gaussian,
$R_k$ is distributed as $\chi^2_{S-1}$ whatever the fault magnitude. The set of
faults with $R_k$ below the $1-\alpha$ quantile $q$ of that distribution is
therefore a confidence set for the cause, with coverage exactly $1-\alpha$. It
is called the certificate. Its size depends on the angles between the
signatures: two faults whose signatures are nearly parallel cannot be
separated, and the companion paper shows that a separation floor of ten
degrees corresponds to a whitened signal-to-noise ratio $\rho = 28.1$.

The certificate is computed from the twin's signatures and evaluated against
observations produced by the process. Its coverage is exact only when the two
agree. This paper asks how far they may disagree.

\subsection{Three classes of parameter}
\label{sec:classes}

The quantities that define the twin enter it in different places, and on that
basis they fall into three classes (Table~\ref{tab:classes}).

\emph{Prediction-only} quantities change a healthy build and cannot change any
fault signature. The keyhole coefficient $C$ is the one example: it corrects
the predicted depth, and the instruments read only the surface.

\emph{Diagnosis-only} quantities change fault signatures and cannot change a
healthy build. The coupling constants that carry each fault into the model act
only when a fault is present. The noise model acts only through whitening.
A physical channel that the twin omits, such as beam scattering by the vapour
plume, acts only when the gas fault is present. The numerical resolution of the
instrument model acts only on the instrument readings, while the width and
depth come from analytic expressions.

\emph{Shared} quantities change both: absorptivity, conductivity, heat capacity,
and the power and speed at which the twin believes the machine operates.

\begin{table}[t]
\centering
\caption{The three classes of twin quantity. Validation against healthy builds
is sensitive to the first and third rows and, by construction, cannot detect an
error in the second.}
\label{tab:classes}
\small
\begin{tabular}{@{}lll@{}}
\toprule
class & quantities & visible to validation \\
\midrule
prediction only & keyhole coefficient & yes \\
diagnosis only & fault couplings, noise model, scattering channel, & no \\
 & \quad instrument-model resolution & \\
shared & absorptivity, conductivity, heat capacity, & yes \\
 & \quad operating power and speed & \\
\bottomrule
\end{tabular}
\end{table}

The partition is a property of the model's structure, not of any numerical
result, and it is checked in code: changing any diagnosis-only quantity leaves
every predicted width and depth identical to the last bit, and changing the
keyhole coefficient leaves the Jacobian identical to the last bit. It already
answers part of the question. Validation against healthy builds can say nothing
about a diagnosis-only quantity, however much data it uses, and it spends part
of its effort on $C$, which diagnosis does not use.

%% ===========================================================================
\section{What the certificate requires of the twin}
\label{sec:theory}

This section states six results. Each is proved in the text or in
Appendix~\ref{app:proofs} and checked by simulation in the accompanying code,
which reports sixteen checks, all passed. Throughout, $\rho$ is the true
whitened signal norm, $q = q_{1-\alpha}(\chi^2_{S-1})$ is the test threshold,
and the angle between two lines is the folded angle
$d(u,v) = \arccos|\langle u, v\rangle| \in [0^\circ, 90^\circ]$.

\subsection{Only directions matter}

\begin{theorem}[Magnitude invariance]
\label{thm:magnitude}
The residual \eqref{eq:residual} depends on the twin's signature for fault $k$
only through its direction. Multiplying the signature by any nonzero constant,
positive or negative, leaves every residual, every certificate and every
coverage probability unchanged.
\end{theorem}

\begin{proof}
The signature enters \eqref{eq:residual} only through the unit vector $u_k$,
and only inside a square, so neither its length nor its sign can affect $R_k$.
\end{proof}

The consequence for the twin is that an error in the size of a fault's effect,
however large, costs the certificate nothing. Among the coupling constants this
separates two cases. The geometry fault reaches the model through a single
preheat constant, so an error in that constant scales the geometry signature and
is free. The layer fault reaches two physical channels through two constants, so
an error in their ratio rotates the layer signature and is not free.

\subsection{What a directional error costs}

\begin{theorem}[Coverage under directional error]
\label{thm:coverage}
Let the twin's direction for the true cause differ from the true direction by
the angle $\delta$, let the twin's noise model be correct, and let the
observation be $z_w = \rho u + e$ with $e \sim N(0, I_S)$. Then the twin's
residual for the true cause is distributed as the noncentral chi-squared
$\chi^2_{S-1}(\rho^2 \sin^2\delta)$, and the nominal $1-\alpha$ certificate
covers the true cause with probability
\begin{equation}
P\bigl(\chi^2_{S-1}(\rho^2\sin^2\delta) \le q\bigr).
\label{eq:coverage}
\end{equation}
\end{theorem}

\begin{proof}
Let $v$ be the twin's direction. The residual is $\|P_v^\perp z_w\|^2$ with
$P_v^\perp = I - vv^\top$. The noise term $P_v^\perp e$ is standard normal on the
$(S-1)$-dimensional complement of $v$. The signal term $\rho P_v^\perp u$ is a
fixed vector in that complement with squared length $\rho^2\sin^2\delta$. The
squared norm of a standard normal vector plus a fixed offset is noncentral
chi-squared with the squared length of the offset as noncentrality.
\end{proof}

\begin{corollary}[Stronger signals require a more accurate twin]
\label{cor:budget}
Let $\lambda^*$ be the noncentrality at which \eqref{eq:coverage} falls to a
target coverage. The largest directional error that keeps coverage at the
target is
\begin{equation}
\delta_{\max}(\rho) = \arcsin\bigl(\sqrt{\lambda^*}/\rho\bigr),
\label{eq:budget}
\end{equation}
which decreases as $1/\rho$.
\end{corollary}

For three instruments, a nominal 95\,\% certificate and a target of 90\,\%,
$\lambda^* = 0.624$, and the budget is $3.24^\circ$ at $\rho = 14$,
$1.62^\circ$ at $\rho = 28$, $0.81^\circ$ at $\rho = 56$ and $0.38^\circ$ at
$\rho = 120$. At the signal strength that the ten-degree floor implies, a twin
must reproduce every fault's direction to within $1.6^\circ$.

The corollary has a consequence that runs against intuition. A strong fault is
easy to detect, and one might expect it to be easy to attribute with an
imperfect twin. The opposite holds: the stronger the signal, the more of it
leaks out of the twin's slightly wrong direction into the residual, and the
sooner the true cause is rejected. A twin adequate for weak faults can be
inadequate for strong faults of the same type.

\subsection{What a directional error does to the reported geometry}

\begin{proposition}[Verdict stability]
\label{prop:verdict}
If the twin's directions for faults $k$ and $l$ are in error by $\delta_k$ and
$\delta_l$, the angle the twin reports between them differs from the true angle
by at most $\delta_k + \delta_l$, and the bound is attained.
\end{proposition}

\begin{proof}
The folded angle is a metric on lines through the origin. The triangle
inequality applied twice gives the bound. It is attained when all four
directions lie in one plane and the errors are in opposite senses.
\end{proof}

A separability verdict for a pair, whether its angle exceeds the floor, is
therefore guaranteed to agree between twin and truth whenever the true angle
lies further than $\delta_k + \delta_l$ from the floor.

\subsection{Errors in the noise model}

\begin{proposition}[Noise-model error]
\label{prop:noise}
Let the twin whiten with noise levels $\sigma_{\text{twin}}$ while the
instruments have $\sigma_{\text{true}}$, and let
$D = \mathrm{diag}(\sigma_{\text{true}}/\sigma_{\text{twin}})$. If the twin's
Jacobian is correct, its residual for the true cause is the quadratic form
$e^\top D P^\perp D e$ with $P^\perp$ the projection off $Du$. Its distribution
is a weighted sum of $S-1$ independent $\chi^2_1$ variables and does not
depend on $\rho$.
\end{proposition}

\begin{proof}
The twin's observation is $D(\rho u + e)$ and its direction for the true cause
is $Du/\|Du\|$, which is parallel to the signal. The signal is therefore
removed exactly by the projection, and only the noise remains.
\end{proof}

A common error $D = cI$ gives coverage $F_{\chi^2_{S-1}}(q/c^2)$. Assuming the
noise to be 10\,\% lower than it is gives 91.6\,\% coverage, 25\,\% lower gives
85.3\,\%, and 50\,\% lower gives 73.6\,\%. Assuming it higher is conservative.
Errors in the ratios between instruments are evaluated by Imhof's method
\cite{Imhof1961}. The contrast with Theorem~\ref{thm:coverage} is the useful
part of the result: an error in the fault model costs more as the signal grows,
while an error in the noise model costs a fixed amount.

\subsection{Build context}

The certificate paper shows that regressing the instrument readings on build
context that the machine already records, such as the layer number and the
local geometry, raises the unique-attribution rate from 3.5\,\% to 95.3\,\%. In
that regime the signature of fault $k$ is the Kronecker product
$u_k \otimes b_k$ of its instantaneous direction and its whitened loading on the
build plan, and the whitener is itself a Kronecker product.

\begin{proposition}[Build context carries the error unchanged]
\label{prop:kron}
If the twin's instantaneous direction is in error by $\delta_J$ and its profile
loading by $\delta_B$, its build-context direction is in error by
$\delta_{\text{big}}$ with $\cos\delta_{\text{big}} = \cos\delta_J\cos\delta_B$.
With the profiles right, $\delta_{\text{big}} = \delta_J$ exactly.
\end{proposition}

\begin{proof}
For unit vectors,
$\langle u\otimes b, v\otimes b'\rangle = \langle u,v\rangle\langle b,b'\rangle$.
\end{proof}

Build context neither amplifies nor attenuates an error in the thermal twin.
It does, however, raise the whitened signal at a fixed fault magnitude, by the
norm of the whitened loading, which is between 9.3 and 11.5 for the five faults
in the build plan used here. By Corollary~\ref{cor:budget} the budget falls in
proportion. At matched fault severity, build context therefore requires a twin
roughly six times more accurate in direction than the instantaneous suite does.
The resolving power and the fidelity requirement have the same source.

\subsection{Measuring a direction instead of modelling it}

Theorem~\ref{thm:magnitude} implies that diagnosis needs from the twin only a
set of directions. For a fault that can be commanded on a healthy machine, such
as a step in power, in speed or in layer thickness, the direction can be
measured instead.

\begin{proposition}[Commanded steps]
\label{prop:commanded}
Command a step of whitened strength $\rho_c$ a total of $N$ times and take the
direction $v$ of the mean response. With $\lambda = \rho_c\sqrt{N}$, the error
of $v$ satisfies exactly
\begin{equation}
\sin^2\hat\delta = \frac{W}{W + (\lambda + z_1)^2},
\qquad W\sim\chi^2_{S-1},\; z_1 \sim N(0,1)\ \text{independent},
\end{equation}
and a later fault of strength $\rho$ is certified with coverage
$E\bigl[P(\chi^2_{S-1}(\rho^2\sin^2\hat\delta)\le q)\bigr]$. For large
$\lambda$ this depends on $\rho/\lambda$ alone, and keeping it at a target
requires $N \ge (\rho/(\kappa^*\rho_c))^2$ for a constant $\kappa^*$ fixed by
$S$ and the target.
\end{proposition}

\begin{proof}
Scaling the mean response by $\sqrt N$ gives $\lambda u + z$ with
$z \sim N(0, I_S)$. The component of $z$ along $u$ is $z_1$, the squared norm of
the rest is $W$, and the stated form follows. The coverage is
Theorem~\ref{thm:coverage} averaged over $\hat\delta$. As $\lambda$ grows,
$\rho^2\sin^2\hat\delta \to (\rho/\lambda)^2 W$.
\end{proof}

For three instruments and a 90\,\% target, $\kappa^* = 0.549$. When the
commanded step is as strong as the fault to be diagnosed, four repetitions
suffice, at every signal strength from $\rho = 14$ to $\rho = 120$. A modelled
direction has a fixed error whose cost grows with $\rho$; a measured direction
has an error that shrinks with the same signal that makes it expensive.

%% ===========================================================================
\section{Twin experiments}
\label{sec:experiments}

%%EXPERIMENTS%%

%% ===========================================================================
\section{Results}
\label{sec:results}

%%RESULTS%%

%% ===========================================================================
\section{Discussion}
\label{sec:discussion}

%%DISCUSSION%%

%% ===========================================================================
\section{Conclusions}
\label{sec:conclusions}

%%CONCLUSIONS%%

\section*{Reproducibility}
All code is in the repository directory \texttt{cirfcms2027/fidelity}.
\texttt{verify.py} checks the six results by simulation; \texttt{study.py}
scores every twin; \texttt{results.py} produces the numbers quoted here, and
\texttt{claims.py} checks each quoted number against them.

\section*{Declaration of generative AI and AI-assisted technologies}
During the preparation of this work the authors used Claude Code in order to
perform formal analysis and to write the original draft. After using this
service the authors reviewed and edited the content as needed and take full
responsibility for the content of the published article.

\bibliographystyle{elsarticle-num}
\bibliography{refs}

\appendix
\section{Proofs}
\label{app:proofs}
%%PROOFS%%

\end{document}
